\PassOptionsToPackage{hypertexnames=false}{hyperref}
\documentclass[twocolumn,tighten,twocolappendix]{aastex701}

\usepackage{newtxtext,newtxmath}
\usepackage{needspace}
\usepackage{orcidlink}
% \usepackage{microtype}
\usepackage{enumitem}
\usepackage{mathtools}

\usepackage{fancyhdr}
\lhead{\footnotesize \textsc{The Astrophysical Journal}, 1003:102 (18pp), 2026 May 20}
\rhead{\footnotesize \href{https://doi.org/10.3847/1538-4357/ae60fd}{https://doi.org/10.3847/1538-4357/ae60fd}}
\cfoot{\thepage}
\pagestyle{fancy}
\renewcommand{\headrulewidth}{0pt}
\setlength{\topmargin}{0.25in}
\setlength{\headsep}{0.125in}
\setlength{\textheight}{60\baselineskip}
\setlength{\footskip}{20pt}

\journalinfo{The Astrophysical Journal, 1003:102 (18pp), 2026 May 20}

% \usepackage{caption}
% \captionsetup{font={stretch=0.5}}

\renewcommand{\sectionautorefname}{Section}
\renewcommand{\appendixautorefname}{}
\renewcommand{\subsectionautorefname}{Section}

\newcommand{\zacc}[1]{\ensuremath{z_\mathrm{acc}}}
\newcommand{\ftr}{\ensuremath{w_\mathrm{tr}}}

\newcommand{\subfind}{\textsc{Subfind}}
\newcommand{\sublink}{\textsc{Sublink}}

\newcommand{\Msun}{\ensuremath{\mathrm{M}_\odot}}
\newcommand{\Mdyn}{\ensuremath{M_\mathrm{dyn}}}
\newcommand{\Mstar}{\ensuremath{M_\star}}
\newcommand{\kms}{\ensuremath{\text{km}~\text{s}^{-1}}}
\newcommand{\osim}{\mathord{\sim}}

\begin{document}

\title{\vspace{-\baselineskip}\Large How to Build an Empirical Speed Distribution\\ for Dark Matter in the Solar Neighborhood}
\shorttitle{How to Build an Empirical Speed Distribution}

\author{\rm Tal Shpigel~\orcidlink{0009-0003-5629-5848}}
\affiliation{\rm Department of Physics, Princeton University, Princeton, NJ 08544, USA; \href{mailto:talshpigel@princeton.edu}{talshpigel@princeton.edu}, \href{mailto:dfolsom@princeton.edu}{dfolsom@princeton.edu}}
\email{talshpigel@princeton.edu}

\author{\rm Dylan Folsom~\orcidlink{0000-0002-1544-1381}}
\affiliation{\rm Department of Physics, Princeton University, Princeton, NJ 08544, USA; \href{mailto:talshpigel@princeton.edu}{talshpigel@princeton.edu}, \href{mailto:dfolsom@princeton.edu}{dfolsom@princeton.edu}}
\email{dfolsom@princeton.edu}

\author{\rm Mariangela Lisanti~\orcidlink{0000-0002-8495-8659}}
\affiliation{\rm Department of Physics, Princeton University, Princeton, NJ 08544, USA; \href{mailto:talshpigel@princeton.edu}{talshpigel@princeton.edu}, \href{mailto:dfolsom@princeton.edu}{dfolsom@princeton.edu}}
\affiliation{\rm Center for Computational Astrophysics, Flatiron Institute, New York, NY 10010, USA}
\email{mlisanti@princeton.edu}

\author{\rm Lina Necib~\orcidlink{0000-0003-2806-1414}}
\affiliation{\rm Physics Department and Kavli Institute for Astrophysics and Space Research, Massachusetts Institute of Technology, Cambridge, MA 02139, USA}
\email{lnecib@mit.edu}

\author{\rm Mark Vogelsberger~\orcidlink{0000-0001-8593-7692}}
\affiliation{\rm Physics Department and Kavli Institute for Astrophysics and Space Research, Massachusetts Institute of Technology, Cambridge, MA 02139, USA}
\affiliation{\rm The NSF AI Institute for Artificial Intelligence and Fundamental Interactions, Massachusetts Institute of Technology, Cambridge, MA 02139, USA}
\email{[mvogelsb@mit.edu]}

\author{\rm Lars Hernquist~\orcidlink{0000-0001-6950-1629}}
\affiliation{\rm Center for Astrophysics, Harvard \& Smithsonian, 60 Garden Street, Cambridge, MA 02138, USA\\\it Received 2025 December 17; revised 2026 April 13; accepted 2026 April 14; published 2026 May 19}
\email{lhernquist@cfa.harvard.edu}

\shortauthors{Shpigel et al.}


\begin{abstract}
    \noindent
    The dark matter flux in a direct detection experiment depends on its local speed distribution.  
    %
    This distribution has been inferred from simulations of Milky Way--like galaxies, but such models serve only as proxies, given that no simulation directly captures the detailed evolution of our own Galaxy.  
    %
    This motivates alternative approaches that obtain this distribution directly from observations.  
    %
    In this work, we utilize 98 Milky Way analogs from the TNG50 simulation to develop and validate a procedure for inferring the dark matter speed distribution using the kinematics of nearby stars. 
    %
    We find that the dark matter that originated from old mergers, plus that from recent nonluminous accretions, is well described by a Maxwell--Boltzmann speed distribution centered at the local standard-of-rest velocity.  
    %
    Meanwhile, recently accreted dark matter from massive mergers has speeds that can be traced from the associated stellar debris of these events.  
    %
    The stellar populations systematically underestimate the velocity dispersion of their dark matter counterparts, but a simple kinematic boost brings the two into good alignment.  
    %
    Using the TNG50 host galaxies, we demonstrate that combining these two contributions provides an accurate reconstruction of the local dark matter speeds.  
    %
    As  an application of the procedure to our own Galaxy, we utilize stellar kinematic data from Gaia to quantify how the dark matter remnants from the Milky Way's last major merger impact its speed distribution in the solar neighborhood.

    \vspace{0.5\baselineskip}\noindent\textit{Unified Astronomy Thesaurus concepts:} \href{https://astrothesaurus.org/uat/2178}{Galactic archaeology (2178)}; \href{https://astrothesaurus.org/uat/353}{Dark matter (353)}; \href{https://astrothesaurus.org/uat/575}{Galaxy accretion (575)}; \href{https://astrothesaurus.org/uat/1051}{Milky Way dynamics (1051)}
\end{abstract}

\section{Introduction} \label{sec:intro}
\thispagestyle{fancy}

Characterizing the speed distribution of dark matter~(DM) near the Sun is crucial to interpret the data from direct detection experiments, which search for DM interactions in terrestrial detector targets (see, e.g., \citet{1996PhR...267..195J,2013RvMP...85.1561F} for reviews). Variations in the assumed speed distribution can impact the expected number of interaction events in these experiments, affecting the interpretation of a signal as well as inferred constraints. Knowledge of the local DM speed distribution is essential to reduce this source of uncertainty~\citep{2014arXiv1404.4130D, 2017JPhG...44h4001G}. 

The DM phase-space distribution near the Sun is a product of the Milky Way's~(MW's) evolutionary history. In the current cosmological paradigm, DM halos grow hierarchically by accreting smaller satellite systems that are tidally disrupted by the host galaxy's gravitational field~\citep{1978MNRAS.183..341W,2008MNRAS.391.1685S, 2011ApJ...740..102K}. DM debris from the earliest mergers, whose time since accretion is many times its orbital timescale, should reach a quasi-equilibrium state with a speed distribution described by a Maxwell--Boltzmann distribution~\citep{1986PhRvD..33.3495D}. However, DM that is accreted more recently will not be relaxed and may retain kinematic signatures of the merger from which it originated~\citep{2008Natur.454..735D}. In the MW's stellar population, there is evidence for such unequilibrated kinematic structure; groups of stars with particular chemical abundances and orbital dynamics appear differently in this chemodynamical space from the stars thought to have formed in the MW itself---see, e.g., \cite{2020ARA&A..58..205H} for a review. These ex situ stellar populations are thought to originate from previous galactic accretion events.

One approach to model the local DM speed distribution is to simulate the formation of galaxies like the MW in a cosmological setting. This has been done in both DM-only~\citep{2005MNRAS.361L...1W,2006JCAP...01..014H,2008MNRAS.385..236V,2009MNRAS.395..797V,2010JCAP...10..034G,2010JCAP...02..030K,2013MNRAS.430.1722V,2017IJMPA..3230016B} and fully hydrodynamical~\citep{2010JCAP...02..012L,2010MNRAS.406..922T,2014ApJ...784..161P,2016JCAP...05..024B,2016MNRAS.462..663B,2016JCAP...08..071K,2016ApJ...831...93S,2017IJMPA..3230016B,  2020JCAP...07..036B,2020JHEP...07..081H,2020JCAP...11..016P, 2023MNRAS.524.2606L,2023JCAP...05..012N,2024JCAP...08..022S} contexts. These simulations can be used to characterize the halo-to-halo variance in the local DM speeds by sampling many possible formation histories. Using the largest high-resolution suite of simulated MW-like galaxies to date, \cite{2025arXiv250507924F} showed that the \emph{ensemble} of DM speed distributions is well modeled by a Maxwell--Boltzmann distribution.  However, they found that the distributions for individual galaxies do deviate from this model, especially at the high-speed tails. Where our own Galaxy falls on this spectrum of possibilities remains an open question. 

Current simulation efforts are trying to produce MW-like galaxies whose formation histories are in better alignment with our own Galaxy~\citep[e.g.,][]{2022MNRAS.512.5823M, 2023MNRAS.521..995R, 2024ApJ...971...79B}, but the challenge of this task motivates complementary efforts to obtain the DM speed distribution directly from observations.  This can be done by harnessing the correlation between the kinematics of ex situ stars and the DM that is accreted from the same merger. For instance, \citet{2018PhRvL.120d1102H}, using the \textsc{Eris} hydrodynamic simulation~\citep{2011ApJ...742...76G}, demonstrated that the oldest and most metal-poor halo stars trace the velocities of the relaxed DM, as both would have originated from the earliest mergers. Furthermore, stellar substructure, such as debris flows~\citep{2012PDU.....1..155L}, also exhibits similar kinematic features to their corresponding DM distributions, a finding supported by the work of \citet{2015ApJ...807...14L} using the Via Lactea simulation~\citep{2008Natur.454..735D} and stellar tagging catalog from~\cite{2012ApJ...745..142R}. This correspondence between stellar and DM velocities has been further confirmed by \citet{2019ApJ...883...27N}, both for the relaxed component and present-day debris flows, using two different simulated halos from the FIRE suite~\citep{2018MNRAS.480..800H}.\footnote{Using six halos from the Auriga project~\citep{2017MNRAS.467..179G}, \citet{2019JCAP...06..045B} did not find a correlation between old stars and the \emph{total} DM distribution. This is likely because they did not separate out the DM contributions from more recent mergers when making the comparison~\citep{2018arXiv181204114L}.} 

As detailed in \citet{2019ApJ...883...27N}, an empirical approach to reconstructing the local DM speed distribution requires considering different subpopulations of DM, as defined by merger origin.  Focusing on DM accreted from luminous mergers, they demonstrated that the oldest DM (with accretion redshift $\zacc{} \gtrsim 3$) is well traced by the most metal-poor stars.  DM from later  mergers ($\zacc{} \lesssim 3$) may be traced by intermediate-metallicity stars; this is especially true if the DM is in velocity substructure, such as debris flow, but less so for the most recent accretions that leave behind  streams. In this paper, we build  on the work of \citet{2019ApJ...883...27N} in several key ways.  First, we vastly increase the sample size of MW analogs studied and also use a separate galaxy formation model.  Second, we account for the contribution of DM from nonluminous mergers and diffuse accretion~\citep{2011MNRAS.413.1373W, 2015MNRAS.448.2941S, 2010MNRAS.401.1796A}, which we refer to as ``dark accretion'' throughout.

In particular, we use a sample of 98 MW-like galaxies from the TNG50 simulation~\citep{2019MNRAS.490.3196P,2019MNRAS.490.3234N} and find that the DM in the solar neighborhood can be modeled as two distinct components: a background that follows a Maxwell--Boltzmann distribution peaking at the local standard-of-rest speed and a population originating from recent, high-mass accretion events.  The former is the combination of DM accreted from the oldest luminous mergers, as well as that from dark accretion.  By combining a Maxwell--Boltzmann model with a stellar-informed model for the merger component, we demonstrate how to reconstruct the total DM speed distribution in the solar neighborhood, accounting for the contribution of dark accretion. These findings, validated on simulation data, establish a practical framework for building a complete empirical speed distribution for the local DM. 

As a first application of this technique, we update previous empirical distributions for the local DM~\citep{2019ApJ...874....3N, 2024ApJ...974..167Z}.  Our result accounts for both dark accretion, as well as the Gaia Sausage--Enceladus merger \citep[GSE;][]{2018MNRAS.478..611B,2018Natur.563...85H}, believed to be the MW's last major merger. We find that the GSE-contributed DM is modestly slower than the local standard-of-rest speed. The mode of the resulting speed distribution is 11~\kms{} slower than that of the Maxwell--Boltzmann-only model, resulting in an $\osim20\%$ suppression to the high-speed tail.

This paper is organized as follows. \autoref{sec:methods} introduces the TNG50 simulation, the criteria used to select the sample of MW analogs, the methodology for tracking the origin of DM and stellar particles in the solar neighborhood, and the selection of analyzed mergers. \autoref{sec:results} presents the core results of the study, focusing on the kinematic correlations between DM and stellar components within the selected merger sample, and the reconstruction of the local DM speed distribution. \autoref{sec:discussion} applies our formalism to the MW and discusses the implications for direct detection experiments. Finally, \autoref{sec:conclusions} summarizes the main conclusions of the study and their broader implications. The Appendices provide an expanded discussion of the DM--merger association procedure, the offset in velocity dispersion between stars and DM, the GSE observational data, and other supplementary figures.

\rhead{\footnotesize Shpigel et al.}
\section{Methodology}\label{sec:methods}
\subsection{MW Analogs}

The IllustrisTNG project\footnote{The IllustrisTNG data are available online at \url{https://tng-project.org/}.}~\citep{2019ComAC...6....2N} comprises a suite of magnetohydrodynamic cosmological simulations that utilize the moving-mesh code AREPO~\citep{2010ARA&A..48..391S}. The sample of MW analogs used in this work is taken from the highest-resolution simulation in this suite, TNG50~\citep{2019MNRAS.490.3234N, 2019MNRAS.490.3196P}, which covers a volume of $(51.7~\rm{Mpc})^3$ from early times (redshift $127$) to the present day, taking cosmological parameters as measured by the~\citet{2016A&A...594A..13P}. The mass resolution of a DM~(star) particle in TNG50 is $4.5\times10^5$~($\osim8.5\times10^4)~\Msun$.  Details regarding the IllustrisTNG Project can be found in \citet{2019ComAC...6....2N}; here, we only focus on those elements that are most relevant for this study.

At each saved time step (``snapshot'') of the simulation, structures are found within the volume using a Friends-of-Friends~(FoF) clustering algorithm~\citep{1985ApJ...292..371D} that is run only on the DM particles. Other particle types (gas, stars, and black holes) are associated with the FoF group of the nearest DM particle. The resulting FoF groups typically consist of a central halo and its neighboring substructure, but the groups can include multiple systems if the central halos are close enough. Further, the FoF algorithm does not require that the particles in the resulting groups be gravitationally bound. Therefore, the \subfind{} algorithm~\citep{2001MNRAS.328..726S, 2009MNRAS.399..497D} is used to determine the bound structures (i.e., DM halos) within each FoF group at each snapshot. Since this is done on a snapshot-by-snapshot basis, the same physical object will be given different \subfind{} IDs at each snapshot. The \sublink{} algorithm identifies persistent physical halos from the \subfind{} catalog, tracking when \subfind{} halos merge with each other and are disrupted, resulting in a ``merger tree'' that traces the assembly of the \subfind{} halos. \subfind{} provides a total mass \Mdyn{} for each halo, which includes all particles bound to the halo but excludes mass bound to satellite halos. The contribution to \Mdyn{} that comes from star particles is denoted $M_\star$.

MW analogs are selected from the \subfind{} catalog for the present-day snapshot following the procedure outlined in \citet{2025ApJ...983..119F}. To summarize, MW analogs are defined as those halos that 
\begin{enumerate}
    \item have $M_\star$ within $4\times10^{10}~\Msun$ and $7.3\times10^{10}~\Msun$,
    \item are farther than 500~kpc from any halo with \Mdyn{} larger than that of the candidate MW, and
    \item are farther than 1~Mpc from any halo with $\Mdyn>10^{13}~\Msun$.
\end{enumerate}

The first criterion selects galaxies with stellar masses that are comparable to that of the MW, based on $68\%$ confidence intervals quoted in the literature~\citep{2006MNRAS.372.1149F, 2011MNRAS.414.2446M, 2015ApJ...806...96L, 2016ARA&A..54..529B, 2016ApJ...831...71L, 2017MNRAS.465...76M, 2020MNRAS.494.4291C}. The next criterion ensures that the MW analog is not too close to a larger galaxy, excluding interacting galaxies residing in close pairs or groups. This does still allow for an M31-like partner, as it is 760~kpc away~\citep{2017MNRAS.468.3428P, 2018MNRAS.475.4043K,2021ApJ...920...84L, 2022ApJ...928L...5B, 2023PhRvD.107j3003V, 2023MNRAS.521.4863S, 2023ApJ...956...15L}. The last criterion removes galaxies that reside near galaxy clusters, as the Local Group is characterized by its relative isolation and the absence of large galaxy groups in its proximity~\citep{2005AJ....129..178K}. Applying these selection criteria yields a final sample of 98 MW-like galaxies. The MW analogs in this sample have peak dynamical masses of $1.16^{+0.51}_{-0.28} \times 10^{12}~\Msun$.\footnote{Throughout this work, we quote the median, 16th, and 84th percentiles.}

While these selection criteria permit certain features of the Local Group, such as a Large Magellanic Cloud--like satellite or an M31-like partner halo, these are not requirements of our selection. Further, the $\mathcal{O}(10~\mathrm{Mpc})$-scale environment of the TNG50 simulation is modestly dissimilar from that of the Local Group. For example, the simulation volume contains two Virgo-mass galaxy clusters, with $\Mdyn \gtrsim 10^{14}~\Msun$. The MW itself is $\osim 16$~Mpc away from the Virgo cluster, but 81\% of the MW analogs in TNG50 are within 16~Mpc of one of these large clusters, and three analogs are close enough to the largest cluster to be in its FoF group. Therefore, the results of this analysis must be interpreted with the understanding that the simulated environments do not exactly reproduce such Local Group features, which may bias the formation histories of the MW analogs.

\subsection{Tracking of DM and Star Particles}
\label{subsec:mergers}

For each host galaxy, we define the region of interest~(ROI) to be a cylindrical shell spanning galactocentric radii $r\in[6,10]$~kpc and height $|z|\leq 2$~kpc from the disk midplane, chosen to be centered around the Sun's location in our Galaxy, $R_\odot \sim 8$~kpc~\citep{2021A&A...647A..59G,2024A&A...692A.242G}.\footnote{Note that the simulated galaxies are generically of different sizes from the MW, which is itself anomalously compact; see \citet{2007ApJ...662..322H,2013ApJ...779..115B,2016ApJ...833..220L,2020MNRAS.498.4943B,2025MNRAS.540.3493T}. The resulting speed distributions at 8~kpc are therefore not indicative of the MW's speed distribution \citep{2025arXiv250507924F}.} The goal is to determine the origin of the DM and star simulation particles that are located within this volume in order to reconstruct the merger history near the solar neighborhood of the host galaxy.

The tracking procedure for star particles follows that of \citet{2025ApJ...983..119F}. To review, mergers are first extracted from the \sublink{} trees for each MW analog. Next, each star particle is tracked across all simulation snapshots, recording the subhalos to which it is bound to over time. Ex situ stars are identified as those that form outside of the MW analog and that are bound to a merger for at least one snapshot.\footnote{There is, however, ambiguity in $\osim1\%$ of cases \citep{2025ApJ...983..119F}.} These stars are associated with the mergers to which they are bound for the greatest number of snapshots. The definition of ex situ chosen by \citet{2025ApJ...983..119F} prioritizes the purity of the sample, ensuring that the ex situ stars can be identified with a contributing merger. 

Tracking the origin of DM particles in the ROI carries additional challenges relative to the stars. This is because the DM halo surrounding a satellite galaxy is more extended than its  stellar component, so it is easier to miss a satellite's more loosely bound DM particles or to misidentify DM from the host as belonging to the subhalo. The DM tracking algorithm used here is designed to minimize these effects.  For each DM particle, we find the first snapshot where \subfind{} considers it bound to the host MW.  This is labeled as the accretion redshift, \zacc{}, for that particle. If \zacc{} is within the first 2~Gyr of the simulation's start, we do not reconstruct the particle's history any further. If not, then we consider the 2~Gyr preceding the accretion redshift.\footnote{Because stars form over the course of the simulation, this 2~Gyr time window would induce a bias if the algorithm were to be na\"ively applied to the stellar component. Stars formed within 2~Gyr of a merger's infall may not appear as bound to that merger for long enough and would be incorrectly categorized as unassociated. Therefore, we do not apply this algorithm to stars and instead use the algorithm from \citet{2025ApJ...983..119F} described above.} The DM is associated with a subhalo if it is bound to it for more than 70\% of the snapshots in this window.  Those particles that do not satisfy this criterion remain unassociated. \autoref{appendix:merger_association} discusses this algorithm in more detail, including robustness tests for the choice of parameters.  Varying the lookback time from 0.5 to 3~Gyr and the snapshot fraction from 10\% to 90\% does not significantly affect the results.

\begin{figure*}
    \centering
    \includegraphics{fig1a.png}
    \includegraphics{fig1b.png}
    \caption{(Left) Fraction of DM in the solar annulus (6--10 kpc in cylindrical radius and height $|z|\leq2$ kpc) originating from Traceable mergers vs. the fraction accreted prior to redshift 3 (Old Untraceable DM), across the 98 MW analogs.  Points are colored by the number of Traceable mergers and shaped by the presence~(star markers) or absence~(circle markers) of a GSE-like event in that galaxy's history. Marginal histograms show the distribution of each quantity across the sample, given the number of Traceable mergers. For systems with one Traceable merger, $58^{+12}_{-17}\%$ of the DM in the solar annulus is Old Untraceable, while $18^{+15}_{-5}\%$ is Traceable. The remaining $21^{+11}_{-12}\%$ is contributed after redshift three but lacks a substantial population of luminous tracers (Young Untraceable). (Right) Speed distributions for the three DM components in the solar annulus: Old Untraceable~(yellow), Young Untraceable~(magenta), and Traceable~(blue). The solid line represents the median probability density, while shaded regions indicate the 16th--84th percentile range across the sample. While the Young Untraceable component exhibits greater halo-to-halo variance and is biased toward higher speeds, it typically represents a small fraction of the total Untraceable DM, such that the combined distribution is well described by a Maxwell--Boltzmann profile. The Traceable component can be modeled using the kinematics of the associated stellar debris. The DM components (Old Untraceable, Young Untraceable, and Traceable) are defined in \autoref{subsec:mergers}.}
    \label{fig:1}
\end{figure*}

Once the origin of each simulation star and DM particle in the ROI is identified, we divide the particles into three categories based on infall time and the mass of their associated merger:
\begin{enumerate}[wide=0pt, leftmargin=\parindent]
    \item \textit{Traceable}. DM particles with $\zacc{}<3$ that originated from a ``Traceable'' merger, defined as a merger with an infall halo mass of $\Mdyn \geq 10^9~\Msun$ that contributes more than 10\% of the DM in the ROI.
    \item \textit{Young Untraceable}. DM particles with $\zacc{}<3$ that are not Traceable. This population can arise from late-time dark accretion as well as from luminous mergers that do not contribute a substantial fraction of DM to the local neighborhood.
    \item Old Untraceable. DM particles with $\zacc{}>3$, consisting of early-time dark accretion and early luminous mergers, both of which have had sufficient time to virialize in the innermost regions of the MW halo. Both ``Untraceable'' components may include contributions from low-mass subhalos with insufficiently resolved stellar populations. 
\end{enumerate}
Though we do not do so here, the DM from early luminous mergers may be directly modeled using metal-poor stars, as shown by \citet{2018PhRvL.120d1102H} and \citet{2019ApJ...883...27N}.  

The criteria used to select the Traceable component identify large mergers that contribute debris to the ROI, so that this component may be traced by stars that are stripped from the same events. Indeed, every Traceable merger contributes at least 186 star particles ($\osim 1.6\times 10^7~\Msun$) to the ROI. In practice, the requirement that a merger contributes more than 10\% of the DM in the ROI already restricts the sample to relatively massive systems; even without imposing the $\Mdyn\geq10^9~\Msun$ cut, Traceable mergers have peak masses of $8^{+14}_{-5}\times10^{10}~\Msun$. 


\subsection{Identifying GSE-like Mergers}

For the primary results of this work, we do not require that the TNG50 galaxies have merger histories resembling that of the MW. For illustrative purposes, however, we do identify those MW analogs that have a GSE-like merger. The GSE is thought to be the most recent major merger experienced by the MW, occurring 8--10~Gyr ago and depositing $\osim10^{8}$--$10^{9}~\Msun$ of stars primarily to the inner 30~kpc of the Galaxy~\citep{2018ApJ...862L...1D, 2018Natur.563...85H,2018ApJ...856L..26M, 2019NatAs...3..932G, 2019MNRAS.486..378L, 2019ApJ...874....3N, 2020ApJ...897L..18B,2021MNRAS.508.1489F, 2021MNRAS.502.5686I,2021NatAs...5..640M, 
 2021ApJ...923...92N}.

In our simulations, a merger is classified as GSE-like if it 
\begin{enumerate}
    \item contributes more than 50\% of all ex situ stars in the volume $|z|\in[9,15]$~kpc and
    \item exhibits a stellar velocity anisotropy $\beta>0.5$ for the stars within the volume $|z|\in[9,15]$~kpc,
\end{enumerate}
where the velocity anisotropy is defined in terms of the galactocentric spherical velocity components:
\begin{equation}
    \beta=1-\frac{\sigma_\phi^2+\sigma_\theta^2}{2\sigma_r^2} \, . 
\end{equation}
These selection criteria reflect observations of the GSE stellar debris, which accounts for more than 50\% of the ex situ stars within a height $z\in[9,15]$~kpc of the Galactic plane and exhibits significant radial velocity anisotropy~\citep{2018ApJ...862L...1D, 2018ApJ...856L..26M, 2019MNRAS.486..378L, 2019ApJ...874....3N, 2021MNRAS.502.5686I, 2021ApJ...923...92N}. We choose this spatial footprint because it is a region in the MW within which the GSE's stellar contribution is most significant~\citep{2020ApJ...901...48N}, and it was also used in previous studies by \citet{2019MNRAS.484.4471F}. Also, while our GSE-like mergers tend to be more recently accreted than the GSE itself, with infall times of $6^{+3}_{-2}$~Gyr ago, adding in a requirement that the merger be before a redshift of 1 does not significantly change the results.

\needspace{2\baselineskip}
\section{Results} \label{sec:results}

This section characterizes the local speed distribution of accreted DM in the MW analogs of TNG50. We divide the DM into the three components described in \autoref{sec:methods}: Old Untraceable, Young Untraceable, and Traceable. The analysis proceeds in three parts. First, we show that the Untraceable components together follow a Maxwell--Boltzmann distribution; second, we demonstrate a procedure for modeling the Traceable component by its associated stellar debris; and finally, we combine these results to reconstruct the total local DM speed distribution. 

\autoref{fig:1} summarizes the fractions of the three DM components and their speed distributions across the 98 MW analogs. The left panel shows the fraction of the DM in the ROI contributed by Old Untraceable mergers versus the fraction considered Traceable, with histograms on the top and right edges indicating the marginal distributions of these values, split up by the number of Traceable mergers. The markers in the central panel are colored according to this number and are shaped according to the presence~(star marker) or absence~(circle marker) of a GSE-like event in the analog's history. Across the full sample, 70 out of 98 analogs (71\%) host one or more Traceable mergers. Among these 70 systems, 37~(53\%) have exactly one Traceable merger, 28~(40\%) have two, and 5~(7\%) have three, corresponding to an average of 1.1 Traceable mergers per halo across the sample.

The MW itself is characterized by a major merger, the GSE, which occurred $\osim 8$--10~Gyr ago. In our sample, 26 MW analogs have a GSE-like Traceable merger. Of these, the GSE-like merger is usually the only Traceable merger, with only 9 analogs (35\% of GSE-hosting halos) having a second. In the MW analogs with one Traceable merger, most of the DM in the ROI ($58^{+12}_{-17}\%$) is Old Untraceable. $18^{+15}_{-5}$\% of the DM is contributed by its Traceable merger, and the remaining $21^{+11}_{-12}\%$ is Young Untraceable. Requiring that the single Traceable merger be considered GSE-like does not meaningfully impact these fractions. This is expected, since selecting GSE analogs among these halos amounts to a cut on the merger's kinematics rather than on the merger history per se, and it is the merger history that sets the DM composition.

\begin{figure*}
    \centering
    \includegraphics{fig2.png}
    \caption{Stacked speed distributions of the Old Untraceable~(yellow) and Young Untraceable~(magenta) DM components in the solar annulus for three representative MW analogs (each containing one GSE-like Traceable merger) with varying Young Untraceable fractions. The SHM~(black curve) is shown for each halo as a Maxwell--Boltzmann distribution with scale velocity $v_0$ set by the mass enclosed within the solar radius. When the Young Untraceable fraction is small, as in the left panel, the combined distribution is well described by the SHM, yielding an EMD of 9~\kms{} between the SHM and the Untraceable DM. For the middle and right panels, the EMD values increase to 14~\kms{} and 28~\kms{}, respectively, reflecting the growing deviations from a Maxwell--Boltzmann shape as the Young Untraceable fraction becomes larger.} 
    \label{fig:2}
\end{figure*}
\begin{figure}
    \centering
    \includegraphics{fig3.png}
    \caption{EMDs between the SHM and the speed distribution of the Old Untraceable DM~(solid yellow) or the Untraceable components combined~(dashed purple). The Old Untraceable component on its own is well described by a Maxwell--Boltzmann, with an EMD of $13^{+8}_{-6}~\kms{}$, while the Young Untraceable component alone (not shown in the figure) is not, lying $65^{+52}_{-32}~\kms{}$ from the SHM. However, since the Young Untraceable component is typically a small fraction of the DM in the solar annulus, combining the Old and Young Untraceable components together is still well modeled by the SHM, with an EMD of $13^{+13}_{-6}~\kms$.}
    \label{fig:3}
\end{figure}

The right panel of \autoref{fig:1} shows the speed distributions of the three components across all 98 MW analogs. In each bin, the solid line indicates the median value for the probability density $f(v)$ for each of the three components, individually normalized to an integral of unity, and the shaded region indicates the 16th--84th percentile range across the sample. The speed distributions for Traceable DM~(blue), Young Untraceable DM~(magenta), and Old Untraceable DM~(yellow) are shown. The speed  distributions of the Young Untraceable component exhibit large halo-to-halo variance compared to the other two components, and they are also biased toward higher values. However, it typically contributes a small fraction relative to the Old Untraceable contribution, so the speed distribution of the \emph{total} (Old and Young) Untraceable component does not typically exhibit these large non-Maxwellian features.

\subsection{Untraceable Components} \label{subsec:old_untraceable}
Given that the Untraceable components contain contributions from dark accretion, it is difficult to model them directly using stellar tracers. Therefore, we use a simple analytic parameterization to describe this DM distribution. We model the speed distribution of the combined Untraceable population in each halo with a Maxwell--Boltzmann profile, 
\begin{equation}
    f(v) = \operatorname{SHM}(v|v_0) = \frac{4v^2}{\sqrt{\pi}v_0^3}\exp\left(-\frac{v^2}{v_0^2}\right) \, ,
    \label{eq:maxwell}
\end{equation}
which we label as the standard halo model~(SHM).\footnote{While it is common to truncate the SHM at the galactic escape speed,  we do not implement such a truncation for simplicity.} The free parameter of the model, $v_0$, is estimated as
\begin{equation} \label{eq:v0}
    v_0 = \sqrt{\frac{GM(<R_\odot)}{R_\odot}} \, ,
\end{equation}
where $G$ is Newton's constant and $M(<R_\odot)$ is the mass enclosed within a radius $R_\odot=8$~kpc, the midpoint of the ROI. We choose this value to correspond to the predicted circular speed in the halo at the solar radius, assuming a spherically symmetric gravitational potential.

To quantify the agreement between the exact speed distribution and this parameterization, we use the Earth mover's distance~(EMD), 
\begin{equation}
    \operatorname{EMD}(f_1,f_2) = \int\!\left| F_1(v) - F_2(v) \right|\ \mathrm{d}v \,,
\end{equation}
where $f_1$ and $f_2$ are speed distributions and $F_1$ and $F_2$ are their respective cumulative distribution functions. Intuitively, the EMD represents the minimum ``work'' required to transform one probability distribution into the other, where ``work'' is the amount of probability that must be moved times the distance it is moved in speed space. The EMD, therefore, has units of speed. This metric is  indicative of shifts in the mean value of the speed distribution, where the probability density is highest, and should be understood as a rough measurement of the similarity of the distributions. In particular, the EMD is not very sensitive to changes at the high-speed tail, as the work required to shift this portion of the speed distributions is proportional to the probability density there. 

\autoref{fig:2} shows the sum of the Old and Young Untraceable speed distributions for three example MW analogs with increasing fractions of Young Untraceable DM~(3\%, 25\%, and 52\%, respectively).\footnote{Note that the first and last cases are outliers in terms of their Young Untraceable fractions; the second case is more typical for the one Traceable merger population.} The speed distributions for the Old Untraceable DM~(yellow) and Young Untraceable DM~(magenta) populations are stacked such that the top edge of the filled histogram is the distribution for the total Untraceable population. The SHM~(solid black) is shown for each halo as a Maxwell--Boltzmann distribution with characteristic speed $v_0$ set by \autoref{eq:v0}. For the leftmost panel, which has the smallest Young Untraceable fraction, the combined Old and Young Untraceable distribution closely matches the SHM, with an EMD between the two of 9~\kms{}. For the two other panels, with intermediate and high Young Untraceable fractions, the SHM does a poorer job of modeling the distribution, and the EMD increases to 14~\kms{} and 28~\kms{}, respectively. 


\begin{figure*}
    \centering
    \includegraphics{fig4_new.png}
    \caption{Speed distributions of Traceable DM~(filled blue) and stars accreted from the same merger~(solid orange) for the Traceable merger in each of the three MW analogs from \autoref{fig:2}. In each case, the stars underpredict the DM speed. To correct for this, we boost the dispersion of the stellar tracers according to \autoref{eq:boost}, giving the boosted stellar distribution (dashed orange). Before the boost, the EMDs between the DM and stellar speed distributions are 49, 48, and 40~\kms{} for the left, middle, and right panels, respectively. The boosted stellar distribution more closely traces the DM distribution, and the resulting EMDs are reduced to 4, 5, and 5~\kms{}, respectively.}
    \label{fig:4}
\end{figure*}


\autoref{fig:3} shows the EMD values between the SHM and the speed distribution of the Old Untraceable component alone~(solid yellow) and the combined distribution of Old and Young Untraceable components~(dashed purple), across all 98 MW analogs. The Old Untraceable component, which is likely virialized and well phase mixed, has an EMD of $13^{+8}_{-6}~\kms{}$ relative to the SHM, with the largest deviation being 43~\kms{}. The Young Untraceable component alone (not shown in the figure) deviates much more strongly from the SHM parameterization, with an EMD of $65^{+52}_{-32}~\kms{}$, with the most discrepant halo 230~\kms{} from the SHM. When the Untraceable components are taken together, the SHM remains a good model; the distributions are $13^{+13}_{-6}~\kms{}$ from the SHM, with a maximum discrepancy of 41~\kms{}. These distances are similar to the distance between the SHM and the Old Untraceable component alone, indicating that the typically small Young Untraceable fraction in most MW analogs does not significantly affect one's ability to model the combined Untraceable populations of DM with the SHM.

\subsection{Traceable Component} \label{subsec:traceable}

\begin{figure*}
    \centering
    \includegraphics{fig5.png}
    \caption{Correlation between stellar and DM kinematics for the 108 Traceable mergers. The top row shows the mean galactocentric velocities $\bar{v}_i$ for each of the spherical components, $i\in \{r,\phi,\theta\}$, while the bottom row shows the velocity dispersion in each of these components. In each panel, the horizontal axis shows the value for the stars contributed by each merger, while the vertical axis shows the value for the DM contributed by the same merger. Equality is indicated by the dashed gray line, such that probability density along this line indicates good agreement between the DM and the stellar tracers. The probability density across the full sample of mergers is shown in the background~(blue, with darker indicating higher probability), with the GSE-like mergers highlighted using star markers. Across the sample, the DM and stars have similar mean velocities, with a difference in $\bar{v}_i$ of $(-1^{+5}_{-3}, -15^{+16}_{-34}, 0^{+4}_{-6})~\kms{}$ for $(v_r, v_\phi, v_\theta)$. The velocity dispersions, on the other hand, are shifted: the DM exhibits an offset of $\Delta\sigma = 34_{-11}^{+10}~\kms{}$ with respect to the stars, for $\Delta\sigma$ the directionally averaged difference in dispersions (\autoref{eq:delta_sigma}). 
    The dispersions are particularly discrepant ($\Delta \sigma = 43^{+11}_{-10}~\kms{}$) for the GSE-like mergers, which are chosen based on their radially biased stars and have lower stellar tangential dispersions relative to the overall population of mergers. The offset in dispersions biases the stellar tracers to predict lower speeds for the DM, as seen in \autoref{fig:4}, which is corrected for by \autoref{eq:boost}.}
    \label{fig:5}
\end{figure*}


We next consider the DM and stellar debris deposited by Traceable mergers. \autoref{fig:4} compares the speed distributions of Traceable DM~(blue) and the corresponding stars~(solid orange) for the same three MW analogs shown in \autoref{fig:2}. In all three halos, the DM is peaked toward higher speeds than the stars, indicating that the stellar speeds provide a biased model for the Traceable DM distributions. Therefore, we develop a procedure for modifying the stellar speeds to better trace the DM. 

To assess this discrepancy within the full sample of MW-like galaxies, \autoref{fig:5} compares the mean velocities ($\bar{v}_i$, in the top row) and velocity dispersions ($\sigma_i$, in the bottom row) for each spherical component $i\in\{r, \phi, \theta\}$ of the debris deposited by the Traceable mergers. In each panel, the quantity for the Traceable DM is shown on the vertical axis, while the corresponding quantity for the stars is shown on the horizontal axis. The distribution across all 108 Traceable mergers is shown in the background~(blue), with the GSE-like mergers~(star markers) highlighted.

Generally, the mean velocity of the Traceable DM ($\bar{v}^\mathrm{DM}$) is in agreement with the mean velocity for the stars contributed by the same merger ($\bar{v}^\star$), with $\bar{v}^\star - \bar{v}^\mathrm{DM} = (-1^{+5}_{-3}, -15^{+16}_{-34}, 0^{+4}_{-6})~\kms{}$ for the components $(v_r, v_\phi, v_\theta)$. In contrast, the velocity dispersions for the DM ($\sigma^\mathrm{DM}$) are systematically higher than those of the stars ($\sigma^\star$), for which we see $(\sigma^\text{DM} - \sigma^\star) = (27^{+10}_{-11}, 38^{+17}_{-14}, 38^{+14}_{-12})~\kms{}$ for $(\sigma_r, \sigma_\phi, \sigma_\theta)$. This is especially true for the GSE-like mergers, which select for radially biased stellar distributions. For these mergers, the $\sigma_\phi^\star$ tends to be much smaller than the $\sigma_\phi^{\mathrm{DM}}$, even more so than the typical merger.

The systematic offset arises from differences in how quickly the two components are stripped from the infalling merger. The stars, which are deeper in a satellite's potential well, tend to get stripped later in the merging process and are deposited toward the center of the MW. The associated DM, on the other hand, is stripped from the infalling merger sooner and deposited at higher orbital energies. Especially for late-time mergers, when the MW's potential well is deepest and allows for such large differences in orbital energy, this biases the DM speed distribution to higher speeds relative to stars from the same merger, resulting in systematically broader velocity dispersions for the DM. This effect is studied in more detail in \autoref{appendix:boost}.
Due to these systematic discrepancies in velocity dispersions, we implement a boost to the stellar velocities to increase the observed dispersion while maintaining a fixed mean. The resulting boosted velocities, $v^\text{b}$, better trace the DM speed distribution.

For each Traceable merger, we define a one-parameter boosted stellar velocity as follows:
\begin{gather}
    v_i^{\mathrm{b}} = \frac{\Delta\sigma + \sigma_i^\star}{\sigma_i^\star}\left( v_i^\star-\bar{v}_i^\star \right) + \bar{v}_i^\star \, ,\label{eq:boost}\\
    \qquad\text{with} \, \Delta\sigma = \frac{1}{3}\sum_i \left( \sigma^\text{DM}_i - \sigma^\star_i \right)\label{eq:delta_sigma}
\end{gather}
and $i\in\{r,\phi,\theta\}$. This boost keeps the mean value for each component of the stellar velocity fixed, but increases the dispersion by $\Delta\sigma = 34_{-11}^{+10}~\kms{}$, the characteristic offset between DM and stellar velocity dispersions.
The DM and stellar debris from these mergers frequently corotates with the stellar disk, with mean $v_\phi$ displaced from zero toward positive velocities. As a result, simply increasing the dispersions without taking care to maintain a fixed mean would result in poorer agreement between the DM and boosted stars.

The corresponding speed distributions for the boosted stars are shown in \autoref{fig:4}~(dashed orange), and the agreement is evident. While the EMDs between the uncorrected stellar speed distributions and the DM speed distributions are 49, 48, and 40~\kms{} for the left, middle, and right panels, respectively, the boosted stars are only 4, 5, and 5~\kms{} away from the corresponding DM distributions, respectively.

\autoref{fig:6} shows the result of this boost across the sample of 108 Traceable mergers. The EMD between the DM speed distributions and the uncorrected stars~(solid orange) is $47^{+19}_{-15}~\kms{}$, with a maximum EMD of 111~\kms{}. The EMD between the DM speed distributions and the boosted stars (dashed orange) is shown, and the agreement is much better, as the EMDs are reduced to $10^{+11}_{-5}~\kms{}$ with a maximum of 55~\kms{}. These results show that the boosted stellar distribution is a good model for the Traceable DM speeds. Indeed, it is an improvement upon the SHM (solid black), which exhibits an EMD of $18^{+16}_{-8}~\kms$ to the Traceable DM speed distribution.


\begin{figure}
    \centering
    \includegraphics{fig6.png}
    \caption{EMDs between the speed distributions of Traceable DM and the stars accreted from the same merger, across all 108 Traceable mergers. The EMD between the DM and the uncorrected stellar distributions~(solid orange) is $47^{+19}_{-15}~\kms{}$. The EMDs, after applying the boost described in \autoref{eq:boost} to the stellar velocities, are shown as Boosted stars~(dashed orange). The boost reduces the EMD to $10^{+11}_{-5}~\kms{}$, indicating that the boosted stars serve as better tracers for the underlying DM speed distribution. The boosted stellar tracer model is also better than the SHM (solid black), which has an EMD of $18^{+16}_{-8}~\kms$ from the DM speed distribution.
    }
    \label{fig:6}
\end{figure}
\begin{figure*}
    \centering
    \includegraphics{fig7.png}
    \caption{Full reconstruction of the local DM speed distribution. (Left) Speed distributions of the Traceable~(blue), Young Untraceable~(magenta), and Old Untraceable~(yellow) components in the solar annulus for an example MW analog, stacked such that the outer envelope is the total exact DM speed distribution. The reconstructions are comprised of a Maxwell--Boltzmann distribution and an empirical distribution derived from accreted stars according to \autoref{eq:model}. The reconstruction that uses the true values for the Traceable DM fraction (\ftr) and the boost applied to the stellar tracers~($\Delta\sigma$, \autoref{eq:boost}) is shown in light dashed green and has an EMD of 10~\kms{} from the exact speed distribution. The solid dark green band indicates the 16th--84th percentile range across all reconstructions, sampling over the uncertainty induced by varying these parameters. Even with this uncertainty, the sampled distributions closely match the exact DM distribution, with an EMD between them of $8^{+3}_{-2}~\kms{}$. The SHM alone (not shown here) aligns well with the reconstructions for this example. (Right) Distribution of EMDs between the exact speed distributions and the reconstructions across all 70 MW analogs with Traceable mergers, for models using the true values of $\Delta\sigma$ and \ftr{} (dashed light green) and those that sample the uncertainty in these parameters (solid dark green). The true-value models are $10^{+6}_{-4}~\kms{}$ away from the exact speed distributions, while accounting for uncertainty in the parameters modestly inflates these EMDs to $11^{+6}_{-4}~\kms{}$. These results indicate that the model provides a good fit, even when accounting for uncertainty in the parameters. We also compare the exact distributions to the SHM alone~(solid black), which yields an EMD of $14^{+15}_{-6}~\kms{}$. This shows that the SHM provides a reasonable overall approximation, but the inference of the speed distribution improves when including a model of the Traceable component.} 
    \label{fig:7}
\end{figure*}

\subsection{Total Local DM} \label{subsec:total}

With an understanding of the Untraceable DM~(\autoref{subsec:old_untraceable}) and the Traceable DM~(\autoref{subsec:traceable}) components, we are equipped to reconstruct the total DM speed distribution in the ROI. We highlight two parameters for this reconstruction that are difficult to obtain observationally: (i)~the boost factor to correct the distribution of stellar tracers, $\Delta\sigma$ (\autoref{eq:boost}), and (ii)~the fraction of DM in the solar neighborhood traced by these stars, \ftr{}. To illustrate the impact of the uncertainties in these parameters, we perform two versions of the analysis, one in which these values are fixed to the truth for each Traceable merger and one in which they are sampled from the distribution of values extracted from TNG50. 

We neglect uncertainty due to the determination of $v_0$~(\autoref{eq:v0}), as it is well constrained for the MW~\citep{2021EPJC...81..907B}. In the case of analogs with multiple Traceable mergers, we further neglect the uncertainty in the relative weight of each and instead divide the fraction \ftr{} according to the total stellar mass of each merger at infall, $m_\star$, such that each Traceable merger is given a weight $\ftr{}\,m_\star/\Mstar$, for $\Mstar$ the sum of all $m_\star$ values.
In general, the inferred relative weight $m_\star/M_\star$ is $14^{+23}_{-6}$~percentage points from the true value. If the relative weight is instead determined by the stellar mass within the ROI from each merger, the inferred relative weight is $19^{+31}_{-11}\%$~points from the true fraction. This suggests that using $m_\star/M_\star$ at infall provides a reasonable approximation for each merger's contribution to \ftr{}. The discrepancy induced by this approximation is modest relative to the uncertainty arising from the dispersion offset between DM and stars (\autoref{subsec:traceable}) and does not dominate the reconstruction. In principle, the stellar mass--halo mass relation and differences in tidal stripping as the mergers orbit their host will change the true value of this fraction, and we refer the interested reader to \citet{2019ApJ...883...27N} for a more careful study of this relative weighting.

The total speed distribution $f_\mathrm{tot}(v)$ is constructed as a weighted sum of the modeled components, 
\begin{equation}\label{eq:model}
    f_\mathrm{tot}(v) = (1-\ftr)\ \operatorname{SHM}(v | v_0) + \ftr\hspace{-1em}\sum_{\#\ \mathrm{mergers}}\frac{m_\star}{\Mstar} f^\mathrm{b}(v) \, .
\end{equation}
The first term corresponds to the Maxwell--Boltzmann distribution of \autoref{eq:maxwell} with $v_0$ set by \autoref{eq:v0}. The second term is a sum of 
the speed distributions $f^\mathrm{b}(v)$ for each Traceable merger, which are in turn constructed as a Gaussian kernel density estimation of the boosted stellar velocities, $v^\mathrm{b}_\star$, defined in \autoref{eq:boost}.

The left panel of \autoref{fig:7} shows this construction on an example halo, ID~550149, the MW analog presented in the central panels of \autoref{fig:2} and \autoref{fig:4}. The full speed  distribution is broken down into the three constituent components: Old Untraceable DM~(yellow), Young Untraceable DM~(magenta), and Traceable DM~(blue), such that the outer envelope of the histogram is the exact speed distribution, summed across all components. The inferred total speed distribution according to \autoref{eq:model}~(dashed green) is shown, with all parameters (\ftr{}, $v_0$, $m_\star$, and $\Delta\sigma$) fixed to their true values. The resulting distribution has an EMD of 10~\kms{} from the exact distribution. However, these parameters are not easily determined in realistic observational settings. As such, we wish to sample over possible values for them. To this end, we construct probability distributions for $\ftr$ and $\Delta\sigma$ from the simulation data. 

From the left panel of \autoref{fig:1}, it is clear that \ftr{} is sensitive to the number of Traceable mergers; in halos with more Traceable mergers, more DM is contributed by Traceable mergers. Therefore, the probability distribution for \ftr{} used in this reconstruction is equal to the distribution observed in the simulation conditioned on the number of Traceable mergers (i.e., the marginal distribution shown at the top edge of the figure). The halo in \autoref{fig:7} has a single Traceable merger, so \ftr{} is sampled from the $18^{+15}_{-5}\%$ range corresponding to this distribution, constructed from the 37 single-merger MW analogs.

As for $\Delta\sigma$, the sampled distribution is the $43_{-10}^{+11}~\kms$ range obtained for the 26 GSE-like mergers, as the Traceable merger in this halo is itself GSE-like. Were the merger not GSE-like, we would use the full sample of 108 mergers, which has  $\Delta\sigma$ in the range $34_{-11}^{+10}~\kms$. Since there are no additional Traceable mergers in this halo, $m_\star/\Mstar = 1$, and the only contribution to \ftr{} comes from this GSE-like merger.

To account for the uncertainty in these two parameters, we consider each of the $37\times 26 = 962$ combinations of \ftr{} and $\Delta\sigma$ and construct an inferred speed distribution from every choice. This is a maximally conservative approach, since it does not assume any correlation between the two. At each point along the horizontal axis in the left panel of \autoref{fig:7}, the median probability density across these distributions (solid green) is shown, and the 16th--84th percentile range (shaded green) is indicated as a band. The EMD for these reconstructions is $8^{+3}_{-2}~\kms{}$, comparable to the 10~\kms{} found using the true values of the parameters.

The right panel of \autoref{fig:7} shows the results of this reconstruction across the subsample of 70 MW analogs with Traceable mergers.\footnote{The 28 analogs with no Traceable mergers have an EMD of $10^{+5}_{-4}~\kms{}$ from their reconstruction, though this reconstruction consists only of a Maxwell--Boltzmann component, with no stellar tracer component.} The EMD between the exact speed distributions and the reconstructions performed with all parameters fixed to their true values (dashed green) is $10^{+6}_{-4}~\kms{}$. For each halo, we also sample over the uncertainties in \ftr{} and $\Delta\sigma$ as we did for ID~550149. As in the worked example, we set the possible values of \ftr{} equal to the values produced by the halos with the same number of Traceable mergers. Further, when modeling a GSE-like merger, the $\Delta\sigma$ is sampled from the values for all GSE-like mergers; otherwise, it is sampled from the full distribution of $\Delta\sigma$ values. For each MW analog, this gives a range of EMDs corresponding to the distance between the exact speed distribution and the $f_\mathrm{tot}(v)$ model. The choice of $\Delta\sigma$ and $\ftr$ does not strongly impact the recovered EMDs; the distribution of recovered EMDs in each galaxy typically has a width of $\osim2$--$4$~\kms{}  (16th--84th percentile). 
Sampling over all reconstructions\footnote{Note that the number of possible reconstructions varies from analog to analog. We have ensured that each analog contributes the same amount of probability density to the Sampled curve in \autoref{fig:7}.} for all 70 MW analogs with traceable mergers yields the Sampled EMD distribution~(solid green) in \autoref{fig:7}, which has a typical EMD of $11^{+6}_{-4}~\kms{}$, providing a modestly worse---though still comparable---reconstruction than the true parameter values. Together, these two distributions indicate that the reconstruction is still well determined when the exact values of \ftr{} and $\Delta\sigma$ are unknown, suggesting that the method is applicable to observational data. 

For comparison, the right panel of \autoref{fig:7} also compares the exact speed distribution to the SHM alone (solid black). The SHM achieves a median EMD of $14^{+15}_{-6}~\kms{}$, which demonstrates that while it provides a reasonable overall approximation to the exact distributions, the  reconstruction does improve once the Traceable components are explicitly modeled. This is true even in halos with one Traceable merger, when the Young Untraceable fraction comprises a nonnegligible fraction of the DM in the ROI. For halos with especially large Young Untraceable fraction, the reconstructions become $\osim 10~\kms{}$ farther from the exact distribution because the Young Untraceable component is particularly non-Maxwellian. However, the uncertainty in the reconstruction parameters $(\Delta\sigma, \ftr{})$ is larger than the variation induced by the Young Untraceable DM; as such, the Sampled EMDs are largely insensitive to the Young Untraceable fraction.


\section{Discussion} \label{sec:discussion}

\begin{figure}
    \centering
    \includegraphics{fig8_new.png}
\caption{Projected local DM speed distribution for the MW in the Galactocentric frame. The SHM with $v_0=238~\kms{}$ is shown~(solid black), as well as high-eccentricity ex situ stars selected from the \citet{2020A&A...636A..75O} catalog as a model for the GSE stellar debris~(solid orange). These GSE stars are boosted using the $\Delta\sigma$ values extracted from GSE-like mergers in TNG50~(dashed orange). Each of these speed distributions has further uncertainty induced by \ftr{}, the relative weighting of the SHM and the stellar tracers, taken from the single Traceable merger halos in the simulation. Summing these two components yields the full reconstructions~(solid green), which lie $6^{+5}_{-3}~\kms{}$ from an SHM-only model. 
    }
    \label{fig:8}
\end{figure}

This section applies the formalism developed in \autoref{sec:methods} and \autoref{sec:results} to our own Galaxy.  \autoref{fig:8} presents the reconstruction of the local DM speed distribution for the MW in the Galactocentric frame. As in \autoref{subsec:old_untraceable}, we describe the Untraceable DM with the SHM~(black), with scale velocity set to the local standard-of-rest speed $v_0=238~\kms{}$ \citep{2004ApJ...616..872R,2016ARA&A..54..529B,2021EPJC...81..907B,2021A&A...647A..59G}. We neglect uncertainty in this parameter, as it is estimated to be only 1.5~\kms{}~\citep{2021EPJC...81..907B}. The shaded gray band indicates uncertainty in the Untraceable DM fraction, described below.

Though many individual mergers are thought to contribute to the MW's ex situ stellar population, by far the largest contribution around the Sun is from the GSE merger \citep{2018Natur.563...85H,2018ApJ...856L..26M, 2020ApJ...901...48N}. Therefore, we take the GSE to be the only Traceable merger in our Galaxy and assume that no other merger contributes more than 10\% of the local DM. To model this component, we select GSE-contributed stars from the catalog of \citet{2020A&A...636A..75O}, which consists of stars from the Gaia DR2 dataset~\citep{2018A&A...616A...1G} flagged by a neural network as likely being ex situ based on their kinematics. 
To isolate the GSE population from this sample, we apply an eccentricity cut of $e>0.7$ and a spatial cut ($r \in [6,10]$~kpc, $|z|\leq 2$~kpc)~\citep{2020ApJ...901...48N}. 
% As discussed in \autoref{appendix:gaia_GSE}, w
We have tested alternative GSE selections and confirmed that this choice yields a stellar speed distribution consistent with the literature~\citep{2019PhRvD..99b3012E,2019ApJ...874....3N,2019MNRAS.486..378L,2020ApJ...903...25N,2020PhRvD.101b3006O,2024ApJ...974..167Z}. While we take this to be a first approximation to the speed distribution of the GSE debris, a more detailed clustering analysis of the \citet{2020A&A...636A..75O} stars in this region would refine this analysis by reducing contamination from, e.g., high-$\alpha$ disk stars, which may contribute a modest ($\osim 25~\kms$) prograde rotation in our GSE sample. See \autoref{appendix:gaia_GSE} for a more detailed discussion.

The resulting speed distribution for the GSE stars is shown in \autoref{fig:8}~(solid orange), where it is normalized by the weight parameter \ftr{} drawn from the distribution of single Traceable merger MW analogs in TNG50 (see \autoref{subsec:total}). The thinner lines indicate the 16th--84th percentiles of these samples. As described in \autoref{subsec:traceable}, we apply a correction factor to these stellar velocities to give a prediction for the GSE-contributed DM. The values of $\Delta\sigma$ are taken from the GSE-like mergers in TNG50. The boosted stellar distribution~(dashed orange), constructed using this boost, is shown with uncertainty from sampling over both $\Delta\sigma$ and \ftr{}.

The total reconstructed speed distribution, formed by combining the SHM and the boosted stellar component, is shown as a green band that samples over both model parameters, where at each point along the horizontal axis, we have shaded the 16th--84th percentiles of the reconstructions. These sampled distributions have an EMD of $6^{+5}_{-3}~\kms{}$ from the SHM, indicating that the model incorporating GSE stars is similar to the Maxwell--Boltzmann-only model. In contrast, constructing the total speed distribution without applying the boost to the stellar velocities yields an EMD of $17^{+11}_{-6}~\kms{}$, as the unboosted GSE stellar speeds are typically much slower than $v_0$. Even after the boost is applied, the GSE stars are generally slower than $v_0$, leading to a suppression of the probability on the high-speed tail from the $(1-\ftr{})$ weighting of the SHM. Due to this effect, the 95th percentile of DM speed is $467^{+5}_{-5}~\kms{}$ for our model, compared to $470~\kms{}$ for the SHM alone.

We can compare our final results to other predictions in the literature for the local DM speed distribution.  For example, \citet{2019ApJ...883...27N} and \citet{2024ApJ...974..167Z} sought to model the DM arising specifically from the luminous mergers, including the GSE, using ex situ stars as tracers. In the Galactocentric frame, we find that the stellar debris from the GSE is slower than the 238~\kms{} local standard-of-rest speed in the MW, with a mode stellar speed of 90~\kms{}---compared to 150~\kms{} and 170~\kms{} for \citet{2019ApJ...883...27N} and \citet{2024ApJ...974..167Z}, respectively. When we use the boosted stellar model, our mode speed increases to $140_{-12}^{+13}$~\kms{} (see \autoref{fig:8}), comparable to the previously published results.

Because we include the contribution of dark accretion in the total DM distribution, we can compare the final result to the SHM$^{++}$ model~\citep{2019PhRvD..99b3012E}, which adds an analytic model for the GSE to the SHM. \cite{2019PhRvD..99b3012E} estimate that the DM fraction in the solar neighborhood due to the GSE is $\osim10\%$--30\%, consistent with the prediction for \ftr{} (13\%--33\%) that we obtain from TNG50. In the Galactocentric frame, their GSE model has larger velocity dispersions than our reconstruction (see \autoref{fig:14}), but its speed distribution is still peaked below the local standard-of-rest speed. We note that the high dispersion of their GSE model results in a total DM distribution in the heliocentric frame that is slightly faster than predicted for the isotropic SHM. In contrast, our total distribution is peaked below the SHM when boosted to the heliocentric frame.

\section{Conclusions} \label{sec:conclusions}
In this paper, we have studied how to empirically build a model for the local distribution of DM speeds. Using a sample of 98 MW analogs from the TNG50 simulation, we demonstrated how to model the contribution of DM from both dark accretion and luminous mergers in the solar neighborhood, defined as the cylindrical shell $r\in[6,10]$ kpc and $|z|\leq 2$~kpc.
Our results demonstrate that the contribution from dark accretion, as well as  early ($\zacc{} > 3$) luminous mergers, can be modeled as a Maxwellian.  Additionally, the stellar debris from recent and significant  luminous mergers can reliably trace the corresponding DM.

This is the first study showing how to properly treat the contribution of dark accretion to the solar neighborhood.  We accounted for the DM that originated from early-accretion events, as well as late-time accretion from smaller mergers. In general, the former has time to virialize and is, per se, well modeled today by a Maxwell--Boltzmann distribution.  The latter can exhibit more variation in its speed distributions.  However, in most cases, it is subdominant to the early-accretion events, such that the Maxwellian continues to be a good model for the total speed distribution of these two contributions, with typical EMDs of $13_{-6}^{+13}~\kms{}$ from the Maxwellian model. This confirms that this ``Untraceable'' component of the local DM can be modeled with this parameterization.

For the remainder of the DM, which comes from large mergers that deposit significant amounts of debris into the ROI, the speed distribution can be reconstructed using the kinematics of the stars from these mergers. However, the stellar debris from these larger mergers is systematically slower than the corresponding ``Traceable'' DM. This is because the DM is stripped from the infalling mergers first and is deposited at higher orbital energies. Despite this discrepancy, the velocities of the DM and stellar populations are strongly correlated. Applying a boost to its velocity dispersion (\autoref{eq:boost}) enables the stellar debris to serve as a good model for the non-Maxwellian component of the DM speed distribution, with an EMD of $10_{-5}^{+11}~\kms{}$ between the DM and boosted stars. 

We combined these results to reconstruct the total local DM speed distribution by adding the stellar-based model for the Traceable component to the Maxwell--Boltzmann model for the Untraceable background. As shown in \autoref{fig:7}, this reconstruction yields strong agreement with the exact distributions across all 70 MW analogs that have stellar tracers, even when accounting for uncertainty in the relative weighting of the Maxwell--Boltzmann and Traceable components of the model. Sampling over this uncertainty produces models with median EMDs of $11_{-4}^{+6}~\kms{}$ from the exact speed distribution. Importantly, the same agreement holds for the subset of analogs that experienced GSE-like mergers, indicating that our conclusions are consistent with expectations for the MW's accretion history.

Finally, in \autoref{sec:discussion}, we applied this approach to the MW itself. Tracers for the GSE DM are taken to be high-eccentricity ex situ stars from the catalog of \citet{2020A&A...636A..75O}. While the GSE stars themselves are slower than the SHM Maxwellian, with a median speed of $180~\kms{}$ (rather than the $\osim260~\kms{}$ expected from the local standard-of-rest velocity), applying our correction factor to the tracers' velocities yields a median speed of $\osim220~\kms{}$ for the GSE-contributed DM, with the exact value set by the correction factor used. As a result, incorporating the empirical speed distribution of the GSE stars only modestly modifies the prediction set by the SHM alone, at least for the DM speed. It also leads to a modest $\osim20\%$ suppression on the high-speed tail compared to the pure Maxwellian. The resulting models for the MW DM speed distribution are available online.\footnote{\url{https://github.com/Tal-Shpigel/stellar-dm-velocity-distributions}} The results of this work motivate improved reconstructions for the velocity distributions of GSE stars in the solar region, which can be used to continue improving the reconstruction of the local DM speed distribution.

This work focuses only on the DM speed distribution, rather than a directionally sensitive velocity distribution. This is sufficient for predicting the scattering rates in an isotropic detector material, but the directional information provided by the stellar tracers may be leveraged by an anisotropic DM detector. Indeed, \autoref{fig:5} suggests that many mergers deposit debris with nonzero azimuthal velocity, contrary to the typical assumption of isotropy. Implications of this will be explored in further work. 

This paper develops a set prescription for how to build an empirical speed distribution, verifying it with a large sample of MW-like systems from TNG50.  These results motivate further study to stress-test the procedure on increasingly large samples of MW-like systems with improved resolution, as well as variation on the subgrid feedback prescriptions.  The latter is likely an important source of systematic uncertainty that should be accounted for when using empirical speed distributions in direct detection analyses.  In terms of the overall methodology, the conclusions of our work are consistent with the prior work of \cite{2019ApJ...883...27N}, which studied two galaxies using FIRE-2 physics~\citep{2016ApJ...827L..23W}, as well as the results in the recent work by \cite{2026arXiv260325783Z}, which expands the sample size to five FIRE-2 MW-like simulations, studying the five mergers in each galaxy that contribute the most stars to the solar annulus. 

\citet{2026arXiv260325783Z} expands upon the prior work in FIRE-2 by considering 
the contribution of dark accretion.  They find---in agreement with our study---that the Untraceable component is well modeled by a Maxwellian distribution, even though the FIRE-2 halos form later than the halos here and have therefore not had as long for their DM to equilibrate. This may be due to the short dynamical timescales in the inner galaxy. 

We find a modest offset between the ex situ stellar velocity dispersions and the Traceable DM counterparts, on the order of a few tens of kilometers per second. This is in contrast with both \cite{2019ApJ...883...27N} and \cite{2026arXiv260325783Z}, who find a tighter correlation between the two distributions for GSE-like mergers. The difference between the FIRE-2 and TNG50 results can be attributed to multiple effects, both numerical and physical, which are challenging to disentangle at this stage.  

On the numerical side, the Latte suite of FIRE-2 has a higher resolution than TNG50, with over an order of magnitude lower mass per simulation particle. Results from \citet{2025arXiv251204157L} suggest that DM speed distributions around our ROI may be resolution-dependent, and it is yet unclear whether TNG50 has sufficient resolution for convergence. Further, \cite{2026arXiv260325783Z} find that the correlation between DM and stellar speed distributions degrades with lower resolution. Additionally, the work in FIRE-2 uses the \textsc{Rockstar} halo-finding algorithm for defining DM halos rather than \subfind{}, which may affect how well DM can be traced through the simulation history. 

Physically, the local environment of the two simulations differs, as the TNG50 cosmological volume holds two large clusters, while the FIRE-2 halos are simulated as zoom-ins with more isolated initial conditions. This denser environment likely biases the TNG50 halos to form earlier than those in FIRE-2, resulting in more concentrated halos at the present day. This concentration is exacerbated by the baryonic feedback physics, as FIRE-2 has a more aggressive prescription for baryonic feedback, which generally yields less dense halos. Both of these will affect tidal disruption of satellites that are orbiting the host systems and the resulting distribution of orbital energies for the DM and stars. 

In general, it is beneficial to vary over these possible systematics by considering many possible simulations. As the sample size of simulated galaxies increases, the connection between a merger's DM and stars will be probed in more depth and with increasing statistical power. Further, simulations performed at higher resolution will be able to identify smaller mergers than those considered here, allowing more of the MW's known mergers to be traced. Advancing hand-in-hand with these computational efforts, future observational surveys of the stellar halo, such as 4MOST~\citep{2019Msngr.175...23H}, WEAVE~\citep{2024MNRAS.530.2688J}, and the LSST~\citep{2019ApJ...873..111I}, will yield more information about the Galaxy's merger history. These efforts will ensure that a detailed empirical model can be borne out to further improve the precision with which one may infer the local DM speed distribution.

\needspace{2\baselineskip}
\section*{Acknowledgments}
We thank Carlos Blanco, Akaxia Cruz, Andreas Thoyas, and Xiuyuan Zhang for useful conversations. M.L. and D.F. are supported by the Department of Energy~(DOE) under Award No. DE-SC0007968. M.L. is also supported by the Simons Investigator in Physics Award. D.F. is additionally supported by the Joseph H. Taylor Graduate Student Fellowship. L.N. is supported by the Sloan Fellowship and the NSF CAREER award 2337864. This work was performed in part at the Aspen Center for Physics, which is supported by National Science Foundation grant PHY-2210452. The computations in this paper were run on the FASRC cluster supported by the FAS Division of Science Research Computing Group at Harvard University. The IllustrisTNG simulations were undertaken with compute time awarded by the Gauss Centre for Supercomputing~(GCS) under GCS Large-Scale Projects GCS-ILLU and GCS-DWAR on the GCS share of the supercomputer Hazel Hen at the High Performance Computing Center Stuttgart~(HLRS), as well as on the machines of the Max Planck Computing and Data Facility~(MPCDF) in Garching, Germany.

This report was prepared as an account of work sponsored by an agency of the United States Government. Neither the United States Government nor any agency thereof, nor any of their employees, makes any warranty, express or implied, or assumes any legal liability or responsibility for the accuracy, completeness, or usefulness of any information, apparatus, product, or process disclosed, or represents that its use would not infringe privately owned rights. Reference herein to any specific commercial product, process, or service by trade name, trademark, manufacturer, or otherwise does not necessarily constitute or imply its endorsement, recommendation, or favoring by the United States Government or any agency thereof. The views and opinions of authors expressed herein do not necessarily state or reflect those of the United States Government or any agency thereof.


\section*{Data availability}
All data used in this work are publicly available at \href{https://tng-project.org/}{https://tng-project.org/}. The inferred MW DM speed distributions presented in \autoref{fig:8} and \autoref{sec:discussion}, as well as a Jupyter Notebook to use these data, are available at \url{https://github.com/Tal-Shpigel/stellar-dm-velocity-distributions} and preserved on Zenodo at \dataset[doi:10.5281/zenodo.19024864]{\doi{10.5281/zenodo.19024864}}.

\section*{ORCID \MakeLowercase{i}Ds}
\noindent Tal Shpigel~\orcidlinkf{0009-0003-5629-5848}\par
\noindent Dylan Folsom~\orcidlinkf{0000-0002-1544-1381}\par
\noindent Mariangela Lisanti~\orcidlinkf{0000-0002-8495-8659}\par
\noindent Lina Necib~\orcidlinkf{0000-0003-2806-1414}\par
\noindent Mark Vogelsberger~\orcidlinkf{0000-0001-8593-7692}\par
\noindent Lars Hernquist~\orcidlinkf{0000-0001-6950-1629}
\onecolumngrid
\clearpage

\renewcommand{\thefigure}{A1}\begin{figure*}
    \centering
    \includegraphics{fig9.png}
    \caption{Spatial distribution of DM particles in MW 550149, focusing on the Traceable merger highlighted in the middle panel of \autoref{fig:4}. This merger's DM particles are shown 2 Gyr before infall (left), at infall (middle), and 2 Gyr after infall (right). All panels display the same pool of 58,541 DM particles that \subfind{} identifies as bound to the merger 2 Gyr before infall, subsampled such that they are more visible in this image. We apply the DM tracking algorithm to divide the particles into three categories: tagged to this merger (blue filled circles), untagged (gray unfilled circles), and tagged to another merger (orange triangles). The host MW is indicated by a point (purple) with its virial radius shown as a circle (dashed purple). Of the selected particles, 96\% are tagged to this merger, 4\% are not tagged to any merger, and $<0.1\%$ are tagged to another merger. This example demonstrates that the tracking procedure robustly associates the majority of DM particles with their correct merger, with very few left untagged or mistagged. Note that this figure assesses the completeness of the algorithm rather than the purity, as particles that are not bound to this merger 2~Gyr before infall are not shown.}
    \label{fig:9}
\end{figure*}

\twocolumngrid

\renewcommand{\thefigure}{A2}\begin{figure}
    \centering
    \includegraphics[width=\columnwidth]{fig10.png}
    \caption{Median purity and completeness across all 98 MW analogs for nine combinations of merger-tracking parameters (see \autoref{subsec:mergers}). Each point shows the median values computed over all mergers with peak mass above $10^9~\Msun{}$. Completeness is defined as the fraction of particles bound to the merger before infall that are tagged to the merger, while purity is the fraction of particles tagged to the merger that are present in the halo before infall. The combinations vary the lookback time window (1--3 Gyr) and the required fraction of snapshots for which a DM particle must be bound to a subhalo~(50\%--90\%); they are labeled by these two parameters. The blue marker with red edge (2.0 Gyr/70\%) denotes the fiducial choice adopted in this work, which maximizes purity while maintaining a completeness above 0.75.}
    \label{fig:10}
\end{figure}


\appendix
\section{Merger Tagging Systematics} \label{appendix:merger_association}


This Appendix presents systematic checks of the DM--merger tagging algorithm introduced in \autoref{subsec:mergers}. Recall that the procedure identifies DM particles that are contributed by a given merger by considering a fixed time interval before the particle is bound to the MW, requiring that the particle remain bound to the subhalo for a sufficient fraction of the snapshots within that time window. Our fiducial choice of parameters is a time window of 2~Gyr, and we require that a particle be bound to its merger for 70\% of the snapshots within this window.

For reference, \autoref{fig:9} illustrates the performance of this fiducial choice for a single merger, showing the positions of a subsample of DM particles at different times relative to the merger's infall time (from left to right: 2 Gyr before infall, at infall, and 2 Gyr after infall), projected onto a coordinate system centered at the location of the MW. We show DM particles that are tagged to the merger~(blue filled circles), particles that are not tagged to any merger~(gray unfilled circles), and particles tagged to a different merger~(orange triangles). In this case, the algorithm correctly tags approximately 96\% of the particles that are bound to the subhalo 2~Gyr before infall, fails to tag about 4\%, and mistags fewer than 0.1\%. This example demonstrates that our fiducial algorithm correctly finds most of the merger debris and, importantly, does so with extremely low contamination. 

To quantify this more systematically, we consider the completeness and the purity of this choice of algorithm parameters. These are assessed  based on the \subfind{} halos 1~Gyr before infall:
\begin{enumerate}[wide=0pt, leftmargin=\parindent]
    \item \textit{Completeness} is the fraction of DM particles bound to the merger halo in any of the three snapshots 1~Gyr before the merger's infall (i.e., the snapshot closest to 1~Gyr before infall, plus or minus one snapshot) that are tagged to the merger by the algorithm.
    \item \textit{Purity} is the fraction of DM particles tagged to a merger that are bound to the merger in any of the three snapshots, 1~Gyr before infall (again, plus or minus one snapshot).
\end{enumerate}

Algorithms with low completeness are more likely to leave the least-bound DM as untagged. These loosely bound DM particles will generally be deposited at the highest orbital energies and have the highest present-day speeds. Therefore, an algorithm with low completeness will likely underestimate a merger's contribution to the high-speed tail. 
On the other hand, algorithms with low purity lead to reconstructions for a merger's DM speed distribution that are contaminated by particles from a different merger, resulting in unpredictable biases to the speed distribution. 
Ideally, the algorithm will have high completeness \emph{and} purity, but optimizing for one often comes at the cost of the other; more selective algorithms will have higher purity, but lower completeness, and vice versa for less restrictive algorithms. Since the low-completeness case introduces a bias that is easier to predict and understand, we prioritize purity, but still aim for high completeness in our choice of tagging algorithm. 

\autoref{fig:10} summarizes how the purity and completeness vary due to the choice of algorithmic parameters. We test nine different combinations of lookback time window (from 0.5 to 3.0 Gyr) and the required fraction of snapshots for which a DM particle must be bound to a subhalo (from 10\% to 90\%); each parameter point is labeled by these values. For each choice of these parameters, we calculate the purity and completeness across all mergers with peak masses above $10^9$~M$_\odot$ and report the medians.

As stated above, there is a trade-off between purity and completeness: relaxing the tagging criteria---by requiring that a particle be bound to a merger for fewer snapshots---raises completeness, since a larger fraction of a merger's particles are recovered. However, this comes at the cost of lower purity, as more particles are incorrectly tagged as belonging to the merger. Conversely, increasing the required bound fraction or extending the lookback time improves purity but excludes some legitimately bound particles, thus lowering completeness. The fiducial choice adopted in this work, 2.0 Gyr/70\%, provides an optimal balance along this trade-off, achieving a median purity of 0.99 and a completeness of 0.79. This combination ensures that nearly all tagged particles are truly bound to the merger, while retaining a sufficient sample size to characterize the merger's kinematics. 

\autoref{fig:11} tests the sensitivity of our reconstruction results to these parameter variations. The left panel shows the distribution of EMDs between the exact DM speed distribution and the sampled reconstruction~(the solid green line in the right panel of \autoref{fig:7}) for each parameter set. While some choices of algorithm parameters result in greater EMDs---indicating worse speed distribution reconstruction---our fiducial choice is among the best-performing algorithms, and the differences between our choice and the other best-performing choices are minimal.

The right panel of \autoref{fig:11} shows the distribution of the boost factor $\Delta\sigma$, as defined in \autoref{eq:delta_sigma}, for each parameter choice. These boosts are consistently biased toward $\osim30~\kms{}$, indicating that the offset is a real physical effect and not the result of our tagging procedure. 

Though these diagnostics are largely insensitive to the choice of parameters used, we caution that \ftr{} does exhibit fairly large variations. The TNG50 halos are rapidly assembling around a redshift of 3, which is 2~Gyr after the start of the simulation. Algorithms with a smaller time window than this threshold can resolve mergers earlier in the simulation history, meaning less DM is classified as Old Untraceable, and more is considered either Traceable or Young Untraceable. This is purely due to a choice in the threshold between ``Old'' and ``Young'' and does not meaningfully impact the physics presented in this work.

These results demonstrate that our reconstruction method is largely insensitive to the specific parameter choice. The fiducial 2.0 Gyr/70\% choice performs comparably to the best alternatives, producing low EMDs and stable boost factors while maintaining high purity and completeness. For these reasons, we adopt this configuration for the analysis presented in the main text.

\onecolumngrid
\clearpage


\renewcommand{\thefigure}{A3}\begin{figure*}
    \centering
    \includegraphics[width=\linewidth]{fig11.png}
    \caption{Comparison of reconstruction performance across the nine merger-tracking parameter sets, labeled by the time window (in Gyr) and the fraction of snapshots that a particle must be bound to be considered tagged (as a percentage). (Left) Distribution of EMDs between the exact DM speed distributions and the sampled--$\Delta\sigma,\ \ftr$ reconstructions (see \autoref{fig:7}) for all Traceable mergers across the 98 MW analogs. While low-purity choices of parameters (i.e., those with smaller time windows or less restrictive fractions) tend to have larger EMDs---and therefore less accurate reconstructions---our fiducial choice is among the most accurate, with a few comparable algorithms. 
    (Right) Distribution of the boost factor $\Delta\sigma$ (\autoref{eq:delta_sigma}) across all Traceable mergers for each parameter set. The results here are qualitatively the same regardless of the choice of DM tagging parameters. There is consistently an offset between the DM and stellar velocity dispersions, suggesting that this offset is a physical result and not due to our tagging procedure.}
    \label{fig:11}
\end{figure*}

\twocolumngrid

\renewcommand{\thefigure}{B1}\begin{figure}
    \centering
    \includegraphics{fig12.png}
    \caption{Gravitational potential per unit mass, $\phi(r)$, at redshift zero across our 98 MW analogs (solid red, with shaded regions indicating 16th--84th percentiles across the analogs). The vertical bands indicate the galactocentric spherical radii that enclose half of the DM debris ($r_{1/2}^\mathrm{DM}$, green) and half of the stellar debris ($r_{1/2}^\star$, blue). The stars are deposited at radii of $6^{+4}_{-3}$~kpc, while the DM is further from the galactic center, at $36^{+19}_{-14}$~kpc. The DM debris typically has higher orbital energy and exhibits a larger velocity dispersion than the stellar debris, yielding the offset shown in \autoref{fig:5}.}
    \label{fig:12}
\end{figure}
\renewcommand{\thefigure}{B2}\begin{figure}
    \centering
    \includegraphics[width=\columnwidth]{fig13.png}
    \caption{Boost factor $\Delta\sigma$ as a function of the infall time for all Traceable mergers (circles). GSE-like mergers are highlighted separately (stars). Mergers with earlier infall times tend to show smaller dispersion offsets. Those with infall times greater than 11~Gyr have a median offset of $9^{+2}_{-5}~\kms{}$, while those with infall times less than 11~Gyr have a median offset of $15^{+6}_{-7}~\kms{}$. This trend is consistent with the mechanism highlighted in \autoref{fig:12}. When mergers fall in earlier, both their DM and stars are deposited at smaller radii because the MW itself was smaller, reducing the difference in potential energies and thus the dispersion offset.}
    \label{fig:13}
\end{figure}
\renewcommand{\thefigure}{C1}\begin{figure*}
    \centering
    \includegraphics{figC1.png}
    \caption{Velocity distributions of Gaia DR2 GSE stars selected with eccentricity $e>0.7$ and restricted to the ROI (black histograms), shown in spherical coordinates: radial ($v_r$, left), azimuthal ($v_\phi$, center), and polar ($v_\theta$, right). For comparison, we show results from previous studies: \citet[][solid yellow]{2019ApJ...874....3N}, \citet[][dashed turquoise]{2019MNRAS.486..378L}, \citet[][dashed-dotted blue]{2019PhRvD..99b3012E}, \citet[][solid orange]{2020ApJ...903...25N}, and \citet[][dotted green]{2020PhRvD.101b3006O}. The distributions of our selected GSE stars are largely consistent with those reported in earlier works.}
    \label{fig:14}
\end{figure*}


\section{DM and Stellar Dispersions}
\label{appendix:boost}
\setcounter{equation}{0}
\renewcommand{\theequation}{B\arabic{equation}}

\autoref{sec:results} showed that the DM associated with a particular merger tends to have a larger velocity dispersion than the corresponding stellar debris, as highlighted in \autoref{fig:5}. This Appendix investigates the origin of this effect. This offset arises because DM is stripped from the merging subhalos first and is left at larger galactocentric radii and higher orbital energies than the stars. The stars, being more tightly bound, remain in the subhalo for longer and are deposited deeper in the MW's potential well, with lower orbital energies. This difference yields DM speed distributions that are biased higher than those of the stars. 

\autoref{fig:12} illustrates this effect. We model the potential well of each host galaxy using its enclosed mass $M(r)$, with 
\begin{equation}
\phi(r) = \frac{GM(r)}{r}
\end{equation}
for each of the 98 MW analogs at redshift zero (solid red, with shaded regions indicating the 16th--84th percentiles across the hosts). For each Traceable merger, we compute the galactocentric spherical radius enclosing half of its stellar debris, $r_{1/2}^\star$~(dashed blue), and half of its DM debris, $r_{1/2}^\mathrm{DM}$~(dashed green), and find that the stars are deposited $6^{+4}_{-3}$~kpc from the galactic center, while the DM is $36^{+19}_{-14}$~kpc from the galactic center. This demonstrates that the stellar debris is deposited at systematically smaller radii than the DM, corresponding to regions deeper in the potential well, with lower orbital energies.

This difference in orbital energies is most distinct for late-infall mergers, when the host potential is particularly deep. Mergers occurring early in a halo's assembly typically exhibit smaller offsets between the DM and stars. This is shown in  \autoref{fig:13}, which plots the boost factor $\Delta\sigma$ for a given merger against its infall time.  The results are shown for all Traceable mergers, and  GSE-like mergers are highlighted with star-shaped markers. Most of the TNG50 halos assemble early and have accreted $\osim60\%$ of their DM by a lookback time of 11~Gyr ago (redshift $\sim 2.5$). The mergers that occur after this time exhibit greater offsets between their DM and stars, with $\Delta\sigma = 15^{+6}_{-7}~\kms{}$. However, very early mergers tend to show smaller offsets, as the hosts' potential wells have not yet reached their full depth. For these mergers, $\Delta\sigma = 9^{+2}_{-5}~\kms{}$. 

This also gives a physical mechanism to explain why we find bigger offsets between a merger's stellar and DM debris compared to \citet{2019ApJ...883...27N} and \citet{2026arXiv260325783Z}. The FIRE-2 halos used in these works assemble more slowly than those in TNG50, lowering $\Delta\sigma$ for more recent mergers and yielding better agreement between the FIRE-2 DM and stars. Furthermore, the baryonic feedback prescription tends to produce shallower potential wells in FIRE-2 than in TNG50 \citep[e.g.,][]{2025arXiv250114868H}, which would further reduce $\Delta\sigma$.


\section{GSE Stellar Sample} \label{appendix:gaia_GSE}

This Appendix describes the data used for the GSE stellar distribution in \autoref{sec:discussion} and \autoref{fig:8}. We provide an account of the source catalog, our subselection of GSE stars, and compare the  results to other selections performed in the literature.

\citet{2020A&A...636A..75O} produced a catalog of accreted stars from the Gaia DR2 dataset \citep{2018A&A...616A...1G}. They train a deep learning classifier on mock Gaia catalogs produced from the FIRE-2 simulation, such that their classifier is able to transfer chemodynamical correlations derived from the simulation to real-world data. This deep learning technique is advantageous, as extracting ex situ stars from a sample containing many in situ MW stars using simple chemodynamical cuts is prone to contamination from the disk and in situ stellar halo \citep{2022MNRAS.510.5119L,2024MNRAS.527.2165C,2024NewAR..9901706D}. The catalog produced by \citet{2020A&A...636A..75O} reproduces known ex situ structures such as the GSE, Helmi streams, and Nyx. This makes it a natural starting point for constructing a sample of GSE stars. 

The neural network classifier outputs a score, $S$, that corresponds to the likelihood that a star is accreted; it ranges from $0$ to $1$ for the least and most likely accreted stars, respectively. We use a likelihood cut of $S>0.75$, which \citet{2020A&A...636A..75O} show most accurately reproduces the ex situ stellar speed distribution. Using the 6D phase-space information of these stars, we compute their eccentricity following the procedure of \citet{2022MNRAS.510.5119L}.\footnote{See \url{https://github.com/jamesmlane/mw-dfs} for an implementation of this method, which we applied to the Gaia~DR2 data.} In particular, we use \textsc{galpy} \citep{2015ApJS..216...29B} to compute actions through the St\"ackel fudge method of \citet{2012MNRAS.426.1324B} and \citet{ 2018PASP..130k4501M}. This approach locally approximates the axisymmetric \textsc{MWPotential2014} potential \citep{2015ApJS..216...29B} as a St\"ackel function, with focal length estimated according to \citet{2012MNRAS.426..128S}. This allows us to compute orbital eccentricities from the Gaia phase-space data.

We consider stars within the solar annulus ROI~$(r \in [6,10]~\mathrm{kpc},\ |z|\leq2~\mathrm{kpc})$ in Galactocentric coordinates. A complete analysis for identifying the GSE stars in this region would require performing a Gaussian mixture analysis, or other clustering method, similar to what was done by \citet{2020ApJ...903...25N} with the same catalog.  While such an analysis is beyond the scope of this work, as a first approximation, we impose an eccentricity cut of $e > 0.7$, in keeping with the recommendation of \citet{2020ApJ...901...48N}. This recommendation is also supported by the findings of \citet{2022MNRAS.510.5119L}, who suggest using an eccentricity cut to achieve a high-completeness selection, which we find suitable given that \citet{2020A&A...636A..75O} already provide a high-purity population of ex situ stars. The resulting distributions of radial, azimuthal, and polar velocities are shown in \autoref{fig:14}, where we compare our selection (black histograms) to others in the literature, namely \citet[][solid yellow]{2019ApJ...874....3N}, \citet[][dashed turquoise]{2019MNRAS.486..378L}, \citet[][dashed-dotted blue]{2019PhRvD..99b3012E}, \citet[][solid orange]{2020ApJ...903...25N}, and \citet[][dotted green]{2020PhRvD.101b3006O}. Our distributions are broadly consistent with these results, though we recover a modest prograde motion to the GSE stars, with $v_\phi = 25^{+50}_{-55}~\kms$. This is possibly due to contamination from the high-$\alpha$ disk, which does contribute stars to the ROI with ex situ likelihood $S > 0.75$ in the \citet{2020A&A...636A..75O} sample. These stars tend to be on circular orbits, so the eccentricity cut removes the majority of them; 62\% of the stars in the ROI are removed by the eccentricity selection, and the majority of these are prograde.

It is important to note that the recommendation of \citet{2020ApJ...901...48N} is not solely an eccentricity cut, but includes cuts to excise the high-$\alpha$ disk, in situ halo, and Sagittarius dwarf~\citep{1994Natur.370..194I}. We neglect these additional cuts because (i) the \citet{2020A&A...636A..75O} sample does not contain the metallicity information required to completely remove the high-$\alpha$ disk and in situ halo components and  (ii) the Sagittarius dwarf is not expected to have deposited many stars within our ROI, as its debris is primarily at Galactocentric radii $r \gtrsim 15$~kpc \citep[e.g.,][]{2020ApJ...900..103J}. We also checked more stringent eccentricity cuts and found that the resulting velocity distributions are similar, though with narrower azimuthal velocity dispersions (i.e., lower $\sigma_\phi$). Because the choice of $e > 0.7$ yields distributions most similar to those recovered elsewhere in the literature, we adopt it as our fiducial case.

The velocity anisotropy of the adopted GSE sample provides an additional diagnostic of its orbital structure. We find a value of $\beta=0.86$, consistent with the highly radial nature of the GSE stars. This result is in good agreement with other studies of the GSE \citep{2018ApJ...862L...1D, 2018ApJ...856L..26M, 2019MNRAS.486..378L, 2019ApJ...874....3N, 2021MNRAS.502.5686I, 2021ApJ...923...92N}.
% \onecolumngrid

\vfill\null
\bibliographystyle{aasjournal}
\bibliography{main}

\end{document}
